% Generato da docs/64-reflex-direttivita-e-cabinet-nominale.md con Pandoc.
% Modificare il Markdown sorgente e rigenerare il file TeX.
% Options for packages loaded elsewhere
\PassOptionsToPackage{unicode}{hyperref}
\PassOptionsToPackage{hyphens}{url}
\PassOptionsToPackage{dvipsnames,svgnames,x11names}{xcolor}
\documentclass[
  italian,
  11pt,
  a4paper,
]{article}
\usepackage{xcolor}
\usepackage[margin=22mm]{geometry}
\usepackage{amsmath,amssymb}
\setcounter{secnumdepth}{-\maxdimen} % remove section numbering
\usepackage{iftex}
\ifPDFTeX
  \usepackage[T1]{fontenc}
  \usepackage[utf8]{inputenc}
  \usepackage{textcomp} % provide euro and other symbols
\else % if luatex or xetex
  \usepackage{unicode-math} % this also loads fontspec
  \defaultfontfeatures{Scale=MatchLowercase}
  \defaultfontfeatures[\rmfamily]{Ligatures=TeX,Scale=1}
\fi
\usepackage{lmodern}
\ifPDFTeX\else
  % xetex/luatex font selection
\fi
% Use upquote if available, for straight quotes in verbatim environments
\IfFileExists{upquote.sty}{\usepackage{upquote}}{}
\IfFileExists{microtype.sty}{% use microtype if available
  \usepackage[]{microtype}
  \UseMicrotypeSet[protrusion]{basicmath} % disable protrusion for tt fonts
}{}
\makeatletter
\@ifundefined{KOMAClassName}{% if non-KOMA class
  \IfFileExists{parskip.sty}{%
    \usepackage{parskip}
  }{% else
    \setlength{\parindent}{0pt}
    \setlength{\parskip}{6pt plus 2pt minus 1pt}}
}{% if KOMA class
  \KOMAoptions{parskip=half}}
\makeatother
\usepackage{longtable,booktabs,array}
\usepackage{caption}
\newcounter{none} % for unnumbered tables
\usepackage{calc} % for calculating minipage widths
% Correct order of tables after \paragraph or \subparagraph
\usepackage{etoolbox}
\makeatletter
\patchcmd\longtable{\par}{\if@noskipsec\mbox{}\fi\par}{}{}
\makeatother
\usepackage{graphicx}
\makeatletter
\newsavebox\pandoc@box
\newcommand*\pandocbounded[1]{% scales image to fit in text height/width
  \sbox\pandoc@box{#1}%
  \Gscale@div\@tempa{\textheight}{\dimexpr\ht\pandoc@box+\dp\pandoc@box\relax}%
  \Gscale@div\@tempb{\linewidth}{\wd\pandoc@box}%
  \ifdim\@tempb\p@<\@tempa\p@\let\@tempa\@tempb\fi% select the smaller of both
  \ifdim\@tempa\p@<\p@\scalebox{\@tempa}{\usebox\pandoc@box}%
  \else\usebox{\pandoc@box}%
  \fi%
}
% Set default figure placement to htbp
\def\fps@figure{htbp}
\makeatother
\ifLuaTeX
\usepackage[shorthands=off]{babel}
\else
\usepackage[shorthands=off]{babel}
\fi
\setlength{\emergencystretch}{3em} % prevent overfull lines
\providecommand{\tightlist}{%
  \setlength{\itemsep}{0pt}\setlength{\parskip}{0pt}}
% Stili A4 per i capitoli tecnici. Il kit Altro si legge dalla macchina.
% Sovrascrivere \AltroFontPath prima di caricare questo file per un altro catalogo.
\providecommand{\AltroFontPath}{D:/design-systems/intrawelt-design/assets/fonts/}
\setmainfont{Altro-Grotesk-Book.ttf}[
  Path=\AltroFontPath,
  BoldFont=Altro-Grotesk-Medium.ttf,
  ItalicFont=Altro-Grotesk-Book.ttf,
  ItalicFeatures={FakeSlant=0.16},
  BoldItalicFont=Altro-Grotesk-Medium.ttf,
  BoldItalicFeatures={FakeSlant=0.16}
]
\newfontfamily\intraweltDisplay{Altro-Serif-Light.ttf}[Path=\AltroFontPath,Scale=1]
\setmonofont{Consolas}[Scale=MatchLowercase]

\definecolor{intraweltViolet}{HTML}{4D11A7}
\definecolor{intraweltInk}{HTML}{0A0A0A}
\definecolor{intraweltRule}{HTML}{D0D4E3}
\AtBeginDocument{\color{intraweltInk}}

\usepackage{titlesec}
\usepackage{listings}
\lstset{
  basicstyle=\small\ttfamily,
  breaklines=true,
  breakatwhitespace=true,
  columns=fullflexible,
  keepspaces=true,
  showstringspaces=false
}
\titleformat{\section}
  {\intraweltDisplay\Large\color{intraweltViolet}}{\thesection}{0.7em}{}
\titleformat{\subsection}
  {\intraweltDisplay\large\color{intraweltViolet}}{\thesubsection}{0.7em}{}
\titleformat{\subsubsection}
  {\intraweltDisplay\normalsize\color{intraweltViolet}}{\thesubsubsection}{0.7em}{}
\titlespacing*{\section}{0pt}{2.1ex plus 0.4ex}{0.75ex}
\titlespacing*{\subsection}{0pt}{1.6ex plus 0.3ex}{0.55ex}
\titlespacing*{\subsubsection}{0pt}{1.2ex plus 0.3ex}{0.4ex}

\makeatletter
\renewcommand{\@maketitle}{%
  \begingroup
    \raggedright
    \color{intraweltViolet}\intraweltDisplay
    \fontsize{29}{31}\selectfont \@title\par
  \endgroup
  \vspace{0.65em}%
  {\color{intraweltRule}\hrule height 0.5pt}%
  \vspace{0.75em}%
}
% Booktabs usa già questo hook per il colore dei filetti.
\newcommand*{\CT@arc@}{\color{intraweltRule}}
\makeatother
\AtBeginDocument{%
  \ifdefined\heavyrulewidth
    \setlength{\heavyrulewidth}{0.55pt}%
    \setlength{\lightrulewidth}{0.45pt}%
    \setlength{\cmidrulewidth}{0.45pt}%
  \fi
  \ifdefined\captionsetup
    \captionsetup[figure]{font=small,labelfont=bf}%
  \fi
}

\setlength{\parskip}{0.34\baselineskip plus 1pt}
\setlength{\parindent}{0pt}
\setlength{\emergencystretch}{2em}
\raggedbottom
\widowpenalty=5000
\clubpenalty=5000

\usepackage{bookmark}
\IfFileExists{xurl.sty}{\usepackage{xurl}}{} % add URL line breaks if available
\urlstyle{same}
% fallback for those not using the hyperref driver hyperxmp:
\makeatletter
\@ifundefined{xmpquote}{\newcommand{\xmpquote}[1]{#1}}{}
\makeatother
\hypersetup{
  pdftitle={Reflex{,} direttività e cabinet: progetto nominale completato},
  pdflang={it-IT},
  colorlinks=true,
  linkcolor={intraweltViolet},
  filecolor={intraweltViolet},
  citecolor={intraweltViolet},
  urlcolor={intraweltViolet},
  pdfcreator={LaTeX via pandoc}}

\title{Reflex, direttività e cabinet: progetto nominale completato}
\author{}
\date{}

\begin{document}
\maketitle

Stato del 25 novembre 2026. Questa pagina completa i confronti e i dettagli geometrici
rimasti aperti nella \href{sistema-autonomo-driver-e-collaudo.pdf}{relazione 63}. I calcoli,
le tavole e gli assiemi sono stati eseguiti; rimangono nominali finché non esistono misure
degli esemplari e del sistema montato. Le verifiche fisiche restano sospese per ADR-055.
Nessun acquisto, taglio o nuovo profilo operativo è autorizzato da questa consegna.

\section{Risultato e scelta di continuità}\label{risultato-e-scelta-di-continuituxe0}

A24 chiusa con SB16/SB26 rimane la base di studio già proposta. La trasformazione LT
nominale più HP35 raggiunge il bersaglio matematico a 60 Hz; il confronto reflex ora
quantifica un'alternativa senza nascondere condotto, occupazione e velocità. C24 dà più
margine al reflex, mentre A24 impone compromessi più stretti. Non viene cambiata
automaticamente la base del progetto sulla sola simulazione.

Il \href{../../research-vault/completamento-20261125/risultati.json}{rapporto numerico}
contiene 108 scenari reflex, 27 scenari di incrocio/offset e 18 ipotesi elettriche. Conserva
anche i condotti che non entrano: 36 scenari rispettano la condizione geometrica rettilinea,
contando tre valori di perdita per geometria. Le figure sono direttamente apribili:
\href{../../build/completamento-20261125-figure/chiusa-reflex.pdf}{chiusa/reflex},
\href{../../build/completamento-20261125-figure/condotto-ingombro-velocita.pdf}{condotto},
\href{../../build/completamento-20261125-figure/direttivita-ideale.pdf}{direttività},
\href{../../build/completamento-20261125-figure/condensatore-e-livelli.pdf}{condensatore e
livelli}.

\section{Il circuito reflex e la sua verifica}\label{il-circuito-reflex-e-la-sua-verifica}

\begin{figure}
\centering
\pandocbounded{\includegraphics[width=.92\linewidth,keepaspectratio,alt={Confronto nominale della cassa chiusa protetta e dei tre accordi reflex.}]{../../build/completamento-20261125-figure/chiusa-reflex.pdf}}
\caption{Confronto nominale della cassa chiusa protetta e dei tre accordi reflex.}
\end{figure}

Il modello usa frequenza normalizzata \(z=jf/f_s\), rapporto \(h=f_b/f_s\), rapporto dei
volumi \(\alpha=V_{as}/V_b\) e resistenza del condotto rappresentata da \(Q_p\). Si risolve
l'ammettenza della cavità; non si divide per un'impedenza infinita nel punto di accordo
ideale.

\[ Y(z)=z+\frac{h^2}{z+h/Q_p},\qquad D(z)=\left(z+\frac{1}{Q_{ts}}+\frac{1}{z}\right)Y(z)+\alpha . \]

\[ H_w=\frac{zY}{D},\qquad H_p=-\frac{zh^2}{(z+h/Q_p)D},\qquad H_{tot}=H_w+H_p=\frac{z^2}{D}. \]

Il segno del condotto è parte del risultato: sotto accordo i contributi possono cancellarsi.
La prova modifica quel segno, osserva il fallimento e verifica il ripristino del sorgente.
Nel caso senza perdite, con \(h=1\), \(\alpha=\sqrt{2}\) e \(Q_{ts}=\sin(\pi/8)\), il
modello coincide con un Butterworth del quarto ordine: modulo \(1/\sqrt{2}\) a \(f_b\) e
pendenza asintotica di 24 dB/ottava. All'accordo ideale la corsa del cono si annulla mentre
il condotto continua a emettere.

Il contesto bibliografico è il lavoro di
\href{https://aes.org/publications/elibrary/elibrary-page/?id=1967}{Small, 1973, verificato
presso AES}. In questa sessione sono stati letti metadati e abstract; non si dichiara una
lettura integrale dell'articolo. Il circuito qui è derivato esplicitamente e provato sui
casi analitici sopra. Sono omessi perdite della cassa, induttanza della bobina, compressione
e comportamento turbolento: i valori \(Q_p=5,10,20\) sono scenari, non misure.

\section{Lunghezza, volume e velocità del
condotto}\label{lunghezza-volume-e-velocituxe0-del-condotto}

\begin{figure}
\centering
\pandocbounded{\includegraphics[width=.92\linewidth,keepaspectratio,alt={Ingombro rettilineo e velocità di picco per i condotti dello studio.}]{../../build/completamento-20261125-figure/condotto-ingombro-velocita.pdf}}
\caption{Ingombro rettilineo e velocità di picco per i condotti dello studio.}
\end{figure}

Il volume utile e la lunghezza si calcolano insieme. Il tubo occupa volume nella sola
porzione interna alla cavità; il tratto che attraversa il pannello non viene sottratto una
seconda volta. La correzione totale delle estremità è assunta pari a 1,46 raggi interni;
parete del tubo 2 mm. Queste ipotesi andranno confrontate con il condotto effettivo.

\[ f_b=\frac{c}{2\pi}\sqrt{\frac{A_i}{V_b L_{eff}}},\qquad L_{eff}=L+1{,}46r,\qquad V_b=V_0-A_o(L-t). \]

Qui \(L\) è la lunghezza fisica complessiva, \(t\) lo spessore del pannello e \(A_i,A_o\) le
aree interna ed esterna del tubo. Si seleziona la soluzione positiva che lascia un volume
utile fisico. La condizione di ingombro richiede il tratto interno più un diametro libero
davanti alla bocca; non include una flangia commerciale o curve non ancora progettate.

Il volume \(V_0\) sottrae 1,5 L di occupazione complessiva ipotizzata, già comprensiva del
rinforzo. Non è il volume misurato di driver, materiale fonoassorbente e cablaggi. Il conto
attuale non somma il rinforzo una seconda volta.

{\def\LTcaptype{none} % do not increment counter
\begin{longtable}[]{@{}llll@{}}
\toprule\noalign{}
Cabinet & Woofer & Volume prima del tubo L & Guadagno massimo LT+HP35 dB \\
\midrule\noalign{}
\endhead
\bottomrule\noalign{}
\endlastfoot
A24 & SB16 & 9,664 & 3,28 \\
B27 & SB15 & 7,237 & 3,68 \\
C24 & SB16 & 12,016 & 2,60 \\
\end{longtable}
}

La tabella seguente estrae quattro casi con \(Q_p=10\). La risposta include HP35 e nessuna
EQ reflex. Velocità ed escursione sono massimi fra 20 e 500 Hz per un livello di passabanda
ipotetico di 90 dB SPL RMS sinusoidali, singolo diffusore in spazio intero a un metro. Non
si deducono SPL reale, margine termico o rumore del condotto.

{\def\LTcaptype{none} % do not increment counter
\begin{longtable}[]{@{}
  >{\raggedright\arraybackslash}p{(\linewidth - 12\tabcolsep) * \real{0.1429}}
  >{\raggedright\arraybackslash}p{(\linewidth - 12\tabcolsep) * \real{0.1429}}
  >{\raggedright\arraybackslash}p{(\linewidth - 12\tabcolsep) * \real{0.1429}}
  >{\raggedright\arraybackslash}p{(\linewidth - 12\tabcolsep) * \real{0.1429}}
  >{\raggedright\arraybackslash}p{(\linewidth - 12\tabcolsep) * \real{0.1429}}
  >{\raggedright\arraybackslash}p{(\linewidth - 12\tabcolsep) * \real{0.1429}}
  >{\raggedright\arraybackslash}p{(\linewidth - 12\tabcolsep) * \real{0.1429}}@{}}
\toprule\noalign{}
\begin{minipage}[b]{\linewidth}\raggedright
Cabinet
\end{minipage} & \begin{minipage}[b]{\linewidth}\raggedright
Accordo Hz
\end{minipage} & \begin{minipage}[b]{\linewidth}\raggedright
Diametro mm
\end{minipage} & \begin{minipage}[b]{\linewidth}\raggedright
Lunghezza totale mm
\end{minipage} & \begin{minipage}[b]{\linewidth}\raggedright
Risposta 60 Hz dB
\end{minipage} & \begin{minipage}[b]{\linewidth}\raggedright
Mach picco
\end{minipage} & \begin{minipage}[b]{\linewidth}\raggedright
Corsa picco mm
\end{minipage} \\
\midrule\noalign{}
\endhead
\bottomrule\noalign{}
\endlastfoot
A24 & 40 & 30 & 115,6 & -2,78 & 0,037 & 2,47 \\
A24 & 50 & 40 & 128,5 & -2,32 & 0,031 & 2,28 \\
B27 & 45 & 30 & 123,8 & -2,28 & 0,049 & 3,30 \\
C24 & 40 & 30 & 88,2 & -1,72 & 0,045 & 2,60 \\
\end{longtable}
}

Il caso C24/40 Hz è interessante per il bersaglio nominale a 60 Hz; non basta a preferirlo
al cabinet chiuso. Il riferimento Mach 0,05 è soltanto una soglia di confronto ipotizzata.
Bocche, posizione del tubo, turbolenza, risonanze longitudinali e compressione possono
cambiare il giudizio. Le lunghezze più grandi restano nel rapporto con stato geometrico
negativo.

\section{Direttività e centri acustici}\label{direttivituxe0-e-centri-acustici}

\begin{figure}
\centering
\pandocbounded{\includegraphics[width=.92\linewidth,keepaspectratio,alt={Direttività verticale ideale con LR4 a 2600 Hz e condensatore da 47 µF; centri acustici coincidenti in profondità.}]{../../build/completamento-20261125-figure/direttivita-ideale.pdf}}
\caption{Direttività verticale ideale con LR4 a 2600 Hz e condensatore da 47 µF; centri
acustici coincidenti in profondità.}
\end{figure}

Per ciascun driver si considera un pistone di area equivalente alla \(S_d\) pubblicata. Il
modello di apertura circolare è il punto di partenza descritto da
\href{https://www.acs.psu.edu/drussell/demos/baffledpiston/baffledpiston.html}{Daniel A.
Russell, Penn State}; non ne vengono riprodotte immagini o animazioni. Il diagramma è
normalizzato e non sostituisce le risposte complesse del driver sul baffle reale.

\[ D(\theta,f)=\frac{2J_1(ka\sin\theta)}{ka\sin\theta},\qquad a=\sqrt{S_d/\pi},\quad k=2\pi f/c . \]

Il valore sull'asse è uno; oltre il primo nullo si conserva anche il segno. Il programma
integra l'apertura con quadratura deterministica, senza installare SciPy. La prova usa il
valore indipendente \(2J_1(1)=0{,}88010117149\) e il primo zero a 3,83170597021.

Alla somma LR4 si applicano il cammino verticale e un offset acustico ipotetico. Nel piano
orizzontale l'interasse verticale non produce lo stesso ritardo.

\[ H=H_{LP}D_w+H_{HP}D_t e^{jk(d\sin\theta+\Delta z\cos\theta)} . \]

Sono esplorati incroci a 2,0/2,6/3,2 kHz e offset -20/0/+20 mm. I due estremi di incrocio
sono scenari matematici, non nuove raccomandazioni di potenza per il tweeter. Senza offset e
senza condensatore, a 2,6 kHz e ±10° il calcolo dà -1,52 dB per SB16/interasse 132 mm e
-1,20 dB per SB15/interasse 119 mm. Con offset +20 mm, A24 arriva a -5,70 dB da un lato: la
geometria da sola non dimostra il criterio fuori asse. Non viene inventata una compensazione
di ritardo da un disegno meccanico.

\section{Condensatore e livelli: che cosa ora è
calcolato}\label{condensatore-e-livelli-che-cosa-ora-uxe8-calcolato}

\begin{figure}
\centering
\pandocbounded{\includegraphics[width=.92\linewidth,keepaspectratio,alt={Effetto del condensatore e richieste elettriche nei guadagni ipotizzati.}]{../../build/completamento-20261125-figure/condensatore-e-livelli.pdf}}
\caption{Effetto del condensatore e richieste elettriche nei guadagni ipotizzati.}
\end{figure}

Il condensatore da 47 µF viene incluso nella somma, con il solo carico resistivo dichiarato
di 3,2 ohm. Il suo polo è 1058,21 Hz; a 2,6 kHz il ramo perde 0,666 dB e anticipa di 22,15°.
La somma dei due rami ideali scende di 0,489 dB sull'asse. Non è corretto aggiungere quel
componente dopo aver congelato fase e livello del crossover.

\[ H_C(j\omega)=\frac{j\omega RC}{1+j\omega RC}. \]

Le ipotesi elettriche partono da +16 dBu dell'interfaccia, preguadagno -12 dB e limite
software di sviluppo Main -6 dB. Per il ramo woofer si conserva il massimo della LT+HP35
nominale e il margine superiore dello shelving di ascolto. I guadagni 20/26/32 dB
appartengono agli scenari, non al Fosi V3 misurato. Le tensioni sono richieste teoriche
prima del clipping.

{\def\LTcaptype{none} % do not increment counter
\begin{longtable}[]{@{}
  >{\raggedright\arraybackslash}p{(\linewidth - 6\tabcolsep) * \real{0.2500}}
  >{\raggedright\arraybackslash}p{(\linewidth - 6\tabcolsep) * \real{0.2500}}
  >{\raggedright\arraybackslash}p{(\linewidth - 6\tabcolsep) * \real{0.2500}}
  >{\raggedright\arraybackslash}p{(\linewidth - 6\tabcolsep) * \real{0.2500}}@{}}
\toprule\noalign{}
\begin{minipage}[b]{\linewidth}\raggedright
Cabinet, modo ascolto
\end{minipage} & \begin{minipage}[b]{\linewidth}\raggedright
Guadagno ipotetico dB
\end{minipage} & \begin{minipage}[b]{\linewidth}\raggedright
Richiesta woofer Vrms
\end{minipage} & \begin{minipage}[b]{\linewidth}\raggedright
Equivalente resistivo 8 ohm W
\end{minipage} \\
\midrule\noalign{}
\endhead
\bottomrule\noalign{}
\endlastfoot
A24 & 20 & 11,30 & 15,97 \\
A24 & 26 & 22,55 & 63,56 \\
A24 & 32 & 44,99 & 253,03 \\
B27 & 26 & 23,62 & 69,72 \\
C24 & 26 & 20,85 & 54,36 \\
\end{longtable}
}

Con alimentazione 32 V, il tetto ideale BTL è \(32/\sqrt{2}=22{,}63\) Vrms prima di ogni
caduta e limite di corrente. I 44,99 Vrms del caso A24/32 dB sono perciò una richiesta
impossibile in quell'ipotesi, non potenza disponibile. Nemmeno 22,55 Vrms autorizza 63,56 W
continui: la vicinanza al limite ideale e la tenuta del driver richiedono taratura. La
potenza IEC del tweeter è dichiarata dal produttore con passa-alto Butterworth a 2600 Hz, 12
dB/ottava; non è una specifica a banda intera.

\section{Cabinet, frontali e distinta}\label{cabinet-frontali-e-distinta}

I disegni ufficiali di
\href{https://sbacoustics.com/wp-content/uploads/2024/01/6in-SB16PFCR25-8.pdf}{SB16},
\href{https://sbacoustics.com/wp-content/uploads/2020/02/5in-SB15NRX2C30-8.pdf}{SB15} e
\href{https://sbacoustics.com/wp-content/uploads/2020/02/SB26STCN-C000-4.pdf}{SB26} sono
stati riletti e ispezionati. Il diametro posteriore del cestello viene distinto dal cerchio
dei fissaggi. Il gioco di un millimetro sul diametro e la profondità di incasso uguale alla
flangia sono scelte nominali da correggere sul processo reale.

{\def\LTcaptype{none} % do not increment counter
\begin{longtable}[]{@{}
  >{\raggedright\arraybackslash}p{(\linewidth - 10\tabcolsep) * \real{0.1667}}
  >{\raggedright\arraybackslash}p{(\linewidth - 10\tabcolsep) * \real{0.1667}}
  >{\raggedright\arraybackslash}p{(\linewidth - 10\tabcolsep) * \real{0.1667}}
  >{\raggedright\arraybackslash}p{(\linewidth - 10\tabcolsep) * \real{0.1667}}
  >{\raggedright\arraybackslash}p{(\linewidth - 10\tabcolsep) * \real{0.1667}}
  >{\raggedright\arraybackslash}p{(\linewidth - 10\tabcolsep) * \real{0.1667}}@{}}
\toprule\noalign{}
\begin{minipage}[b]{\linewidth}\raggedright
Variante
\end{minipage} & \begin{minipage}[b]{\linewidth}\raggedright
Taglio woofer mm
\end{minipage} & \begin{minipage}[b]{\linewidth}\raggedright
Incasso woofer mm × profondità
\end{minipage} & \begin{minipage}[b]{\linewidth}\raggedright
Taglio tweeter mm
\end{minipage} & \begin{minipage}[b]{\linewidth}\raggedright
Incasso tweeter mm × profondità
\end{minipage} & \begin{minipage}[b]{\linewidth}\raggedright
Interasse mm
\end{minipage} \\
\midrule\noalign{}
\endhead
\bottomrule\noalign{}
\endlastfoot
A24/C24 & 145,5 & 177 × 4,1 & 47,15 & 73 × 3 & 132 \\
B27 & 124,9 & 151 × 6,5 & 47,15 & 73 × 3 & 119 \\
\end{longtable}
}

I centri dei quattro fissaggi seguono i cerchi di targa; rimangono croci di costruzione,
perché un foro pilota dipende dal fissaggio scelto. Il ponte nominale fra i due incassi è 7
mm. Il rinforzo è un settimo pannello con finestra, bordo 20 mm, arretramento di 20 mm dal
lato interno del frontale e quota inferiore 201 mm. L'arretramento evita l'inviluppo
conservativo del tweeter. Il modello controlla collisioni fra pannelli e con gli inviluppi
dei driver.

{\def\LTcaptype{none} % do not increment counter
\begin{longtable}[]{@{}llll@{}}
\toprule\noalign{}
Cabinet & Legno per diffusore L & Solidi di legno & Distinta per coppia \\
\midrule\noalign{}
\endhead
\bottomrule\noalign{}
\endlastfoot
A24 & 6,639 & 7 & 14 pezzi, prima di sfridi e prove di fresatura \\
B27 & 5,844 & 7 & 14 pezzi, prima di sfridi e prove di fresatura \\
C24 & 7,455 & 7 & 14 pezzi, prima di sfridi e prove di fresatura \\
\end{longtable}
}

Le giunzioni sono a battuta incollata nel modello, senza una verifica strutturale della
colla o del serraggio. Sigillatura, assorbente, passacavi e terminali non hanno ingombri
esecutivi definitivi: rimangono compresi nell'occupazione nominale da riscontrare. Non si
inventano fresature per connettori ancora da scegliere.

I file della \href{../../build/consegna-completamento-20261125/README.md}{consegna locale}
comprendono FCStd modificabili, STEP, DXF, tavole PDF/SVG e distinta CSV. FreeCAD 1.1.4 ha
verificato riapertura, sette solidi, volumi, assenza di sovrapposizioni e STEP reimportati.
Una variazione della larghezza di 10 mm è stata ricomputata e ripristinata; il raggio
deliberatamente errato ha fatto fallire il confronto dei volumi. CADCraft 0.20.1 ha riletto
dodici profili per variante, verificando area, perimetro, centro e millimetri.

\section{Sequenza replicabile eseguita}\label{sequenza-replicabile-eseguita}

Sulla postazione Windows, dalla radice di questo repository:

\begin{lstlisting}
cd "E:/diy-2way-monitors-home"
python tools/studio_completamento.py
python tools/figure_completamento.py
python tools/cabinet-dettagli-cadcraft.py --out build/completamento-20261125-cadcraft-verificato
python tools/collauda-cabinet-dettagli.py --host studio --out build/completamento-20261125-freecad-verificato --gui
python tools/consegna_documenti.py
python tools/consegna_completamento.py
python -m pytest tools/test_acustica_cabinet.py tools/test_progetto_cabinet.py tools/test_studio_completamento.py tools/test_consegna_cabinet.py tools/test_figure_completamento.py tools/test_consegna_completamento.py -q
python tools/studio_completamento.py --check
python tools/figure_completamento.py --check
python tools/consegna_completamento.py --check
\end{lstlisting}

I comandi di generazione CAD e figure richiedono destinazioni di lavoro nuove. Con risultati
già presenti, si eseguono prima i controlli; per rigenerare si conserva la cartella
precedente e si ricreano i percorsi indicati. Gli SVG pubblicati invariati vengono
riconosciuti, quelli cambiati passano dal generatore con archivio della versione precedente.
Il nome SSH del collaudo è quello della postazione su cui il comando è stato eseguito; su un
clone va passato il proprio alias verificato. FreeCAD riceve soltanto il banco temporaneo;
non vengono modificati servizi audio o installazioni ordinarie e non vengono rimossi file
remoti.

\section{Troubleshooting osservato}\label{troubleshooting-osservato}

\subsection{FreeCAD rifiuta un alias del foglio}\label{freecad-rifiuta-un-alias-del-foglio}

Osservato su Ubuntu il 25 novembre, FreeCAD 1.1.4: gli alias W, H e T vengono rifiutati,
mentre D e i nomi prefissati DimW/DimH/DimT sono accettati. Una prova sul foglio nudo ha
isolato il problema prima di cambiare il modello. La correzione usa il prefisso Dim per
tutte le quote. Il primo rapporto conserva il fallimento; il successivo collaudo dei solidi
e la prova negativa sono riusciti.

\subsection{Il sandbox Windows nega il file temporaneo della
prova}\label{il-sandbox-windows-nega-il-file-temporaneo-della-prova}

Osservato il 25 novembre: la suite del rapporto ha incontrato PermissionError nella cartella
temporanea, prima di produrre il JSON. La stessa suite eseguita con il permesso di
esecuzione appropriato passa, inclusi alterazione del risultato e pulizia. Non è stato
cambiato il codice fisico per aggirare il problema.

\subsection{La vista PNG del CAD ha sfondo
scuro}\label{la-vista-png-del-cad-ha-sfondo-scuro}

L'esportazione PNG predefinita di CADCraft usa il fondo del modello. Le tavole PDF/SVG sono
su carta bianca; la verifica di stampa è stata fatta sul PDF reso con PdfCraft. Le viste
FreeCAD usano impostazioni GUI isolate e colori del design system, senza cambiare le
preferenze ordinarie del proprietario.

\end{document}
